\section{Специфікація стенда: як його відтворити}
\label{sec:dishspec}

Нижче зібрано все, що потрібно, щоб зібрати такий стенд наново: обладнання, петля, кодування входу, читання виходу, зворотний зв'язок і протокол досліду. Кожен параметр позначено джерелом. «Стаття» --- значення взято з тексту роботи~\cite{Kagan2022}. «Вибір» --- у статті його немає, і для відтворення його доведеться задати самому; ми пропонуємо типове значення і кажемо, як його перевірити.

\subsection*{Обладнання і культура}

\begin{center}
\footnotesize
\setlength{\tabcolsep}{4pt}
\renewcommand{\arraystretch}{1.15}
\begin{tabular}{>{\raggedright}p{4.0cm}>{\raggedright}p{6.9cm}>{\raggedright\arraybackslash}p{1.6cm}}
\toprule
Параметр & Значення & Джерело \\
\midrule
чип & матриця MaxOne: $26\,000$ платинових електродів на майданчику $8$~мм$^2$ & стаття \\
запис & до $1024$ каналів одночасно, $20\,000$ відліків на секунду & стаття \\
стимуляція & до $32$ електродів технічно, у досліді --- $8$ незалежних & стаття \\
клітини & кора зародка миші; нейрони людини зі стовбурових клітин (два способи вирощування); клітини HEK293T (не нейрони) для контролю & стаття \\
висів & близько $10^6$ клітин на чип & стаття \\
вік культури на час досліду & миша: від $\sim10$ діб; людина: $14$--$24$ доби або $\sim80$ діб залежно від способу вирощування & стаття \\
\bottomrule
\end{tabular}
\end{center}

\subsection*{Петля і час}

Петля працює вікнами по $10\ms$ ($200$ відліків): за вікно підраховують імпульси на електродах рухових ділянок, гра оновлюється, і видається чергова стимуляція. Затримка від імпульсу в культурі до відповідного стимулу --- близько $5\ms$. Під час гри сигнал кожного електрода проходить фільтр верхніх частот Бесселя 2-го порядку з частотою зрізу $100$~Гц; модуль сигналу згладжується фільтром нижніх частот Бесселя 1-го порядку з частотою $1$~Гц, і поріг імпульсу пропорційний цьому згладженому модулю. Поріг у шість стандартних відхилень фону стаття вказує для окремих вимірювань активності, а не для гри. Щоб стимуляцію не зараховували як відповідь культури, усі сигнали гасять, коли одночасно надходить понад $15$ великих, вищих за $75$~мВ, імпульсів. Усе це --- стаття.

\subsection*{Кодування м'яча}

\begin{center}
\footnotesize
\setlength{\tabcolsep}{4pt}
\renewcommand{\arraystretch}{1.15}
\begin{tabular}{>{\raggedright}p{3.4cm}>{\raggedright}p{7.5cm}>{\raggedright\arraybackslash}p{1.6cm}}
\toprule
Параметр & Значення & Джерело \\
\midrule
чутлива ділянка & $8$ електродів, розташованих так, щоб сусідні електроди означали сусідні положення м'яча & стаття \\
код місця & номер електрода $k=1,\dots,8$ задає положення м'яча відносно центру ракетки & стаття \\
яка координата & вертикальне положення м'яча; для відтворення --- відносно ракетки, $k=\lceil 8(y-p+\tfrac12)\rceil$ при $y-p\in[-\tfrac12,\tfrac12]$ & стаття; формула --- вибір \\
частотний код & частота стимуляції від $4$ до $40$~імп/с несе відстань до ракетки & стаття \\
напрям шкали & $4$~імп/с, коли м'яч біля протилежної стіни, і лінійно до $40$~імп/с, коли він біля стіни ракетки: $f=4+36\,(1-x/L)$, $x$ --- відстань до стіни ракетки, $L$ --- ширина корту & стаття \\
форма імпульсу & короткий прямокутний двофазний: спершу додатна, потім від'ємна фаза & стаття \\
амплітуда & $75$~мВ; обрано так, щоб не змушувати спрацьовувати загальмовані клітини & стаття \\
тривалість фаз & не вказано; взяти стандартне налаштування системи стимуляції і не змінювати його впродовж досліду & вибір \\
\bottomrule
\end{tabular}
\end{center}

\subsection*{Читання ракетки}

У статті дві рухові ділянки: імпульси в першій рухають ракетку вгору, у другій --- вниз; внесок ділянок зважують, щоб урахувати різну щільність клітин і різну активність~\cite{Kagan2022}. Точної формули не наведено. Для відтворення беремо правило~\eqref{eq:dishreadout}, $\Delta p=c\,(n_\uparrow-n_\downarrow)$ за кожне вікно $10\ms$, з вагою, що вирівнює ділянки за їхньою середньою активністю в спокої (вибір): $n_\uparrow$ і $n_\downarrow$ ділять на середню лічбу своєї ділянки за вікно, виміряну перед грою. Коефіцієнт $c$ обираємо так, щоб різниця в одну середню лічбу зсувала ракетку на $1\%$ висоти корту за вікно (вибір); тоді стала перевага однієї ділянки проводить ракетку через увесь корт за $1$~с.

Розташування ділянок на чипі в тексті статті координатами не задано; є лише схема. Для відтворення досить трьох груп електродів, що не перетинаються: вісім чутливих у центрі та по групі записувальних електродів з обох боків, одна --- «вгору», друга --- «вниз» (вибір).

\subsection*{Гра}

Розмірів корту, довжини ракетки і швидкості м'яча в статті не вказано; сказано лише, що м'яч летить зі сталою швидкістю і відбивається від стін. Для відтворення (вибір): корт $1\times1$, ракетка на правій стінці завдовжки $0.2$ (влучання при $|y-p|<0.1$, як у моделі~\texttt{models/dishbrain\_model.py}), м'яч перетинає корт за $1$~с. Промах без жодного відбитого м'яча в розіграші вважають «ейсом»; розіграш, довший за три удари поспіль, --- «довгим» (стаття).

\subsection*{Зворотний зв'язок}

\begin{center}
\footnotesize
\setlength{\tabcolsep}{4pt}
\renewcommand{\arraystretch}{1.15}
\begin{tabular}{>{\raggedright}p{2.6cm}>{\raggedright}p{8.3cm}>{\raggedright\arraybackslash}p{1.6cm}}
\toprule
Подія & Що подається & Джерело \\
\midrule
удар & передбачуваний стимул: усі $8$ чутливих електродів одночасно, $75$~мВ, $100$~імп/с, $100\ms$; звичайна стимуляція м'яча на цей час переривається & стаття \\
промах & непередбачуваний стимул: $150$~мВ, $5$~імп/с, $4$~с, у випадкові моменти на випадково вибрані з $8$ електроди; потім $4$~с без стимуляції, і м'яч стартує у випадковому напрямку & стаття \\
закон випадковості & не вказано; для відтворення --- моменти імпульсів як пуассонівський потік із середньою частотою $5$~імп/с & вибір \\
\bottomrule
\end{tabular}
\end{center}

\subsection*{Протокол і контроль}

Один сеанс триває $20$~хв; порівнюють перші $5$ і останні $15$~хв (стаття). Три міри результату: середня довжина розіграшу, частка ейсів і частка довгих розіграшів. Умови досліду (стаття):
\begin{itemize}
\item \emph{стимул} --- повна петля, як описано вище;
\item \emph{тиша} --- замість стимулу удару чи промаху стільки ж часу без жодної стимуляції, потім м'яч стартує у випадковому напрямку;
\item \emph{без зворотного зв'язку} --- м'яч після промаху просто відбивається й летить далі, стимуляція положення м'яча триває;
\item \emph{спокій} --- гра йде і ракетка рухається за активністю, але стимуляції немає зовсім;
\item \emph{контролі} --- чип лише із середовищем; клітини HEK293T; «культура в комп'ютері», симуляція замість клітин.
\end{itemize}
Обсяг вибірки в основній серії: $9$ культур миші ($101$ сеанс), $11$ культур людини ($138$), $6$ чипів із середовищем ($80$), $20$ культур у спокої ($42$), $3$ симуляції ($38$), разом $399$ сеансів. У серії про зворотний зв'язок --- $486$ сеансів за $3$ дні по $3$ сеанси на день, умови «стимул», «тиша» і «без зворотного зв'язку» ($n=27$, $17$ і $15$). Зростання середньої довжини розіграшу від перших $5$ хвилин до останніх $15$ отримано лише в живих культур. У серії про зворотний зв'язок значуще зростання в часі --- тільки в умові «стимул»; наприкінці сеансу умова «тиша» все ж перевершувала «без зворотного зв'язку», але з меншим ефектом, а стійкого навчання від сеансу до сеансу за три дні не було~\cite{Kagan2022}.

\subsection*{Цикл стенда крок за кроком}

Кожні $10\ms$ комп'ютер стенда робить таке.
\begin{enumerate}
\item Зчитує кількість імпульсів $n_\uparrow$, $n_\downarrow$ у двох рухових ділянках за минуле вікно і нормує їх на середню лічбу ділянки в спокої.
\item Зсуває ракетку на $\Delta p=c\,(n_\uparrow-n_\downarrow)$ і обмежує її краями корту.
\item Зсуває м'яч на один крок; відбиває його від верхньої, нижньої та лівої стінок.
\item Якщо м'яч дійшов до правої стінки: при $|y-p|<0.1$ --- удар, лічба розіграшу зростає на один, подається стимул удару ($100\ms$); інакше --- промах, розіграш записується, подається стимул промаху ($4$~с), потім $4$~с паузи, і м'яч стартує знову.
\item Інакше обирає електрод $k$ за вертикальним положенням м'яча і частоту $f$ за відстанню до стіни ракетки та подає на електрод $k$ двофазні імпульси $75$~мВ із частотою $f$ до наступного вікна.
\end{enumerate}
Цього досить, щоб зібрати стенд, сумісний з опублікованим, і перевірити головне питання розділу: чи навчає культуру передбачуваність, чи проста різниця між відповіддю на удар і на промах. Для цього до трьох умов статті варто додати четверту, якої в ній немає: після промаху --- \emph{передбачуваний} стимул тієї самої тривалості й енергії, що й непередбачуваний. Якщо культура навчається і так, річ у різниці, а не в передбачуваності.
